\documentclass[12pt]{amsart} 
\usepackage{amsfonts}
\usepackage{amscd}
\usepackage{amsmath}
\usepackage{amssymb}
\usepackage{amsthm}
\usepackage{enumitem}
\usepackage{pgf,tikz,tikz-cd}
\usepackage{mathrsfs}
\usepackage[margin=1in]{geometry}
\usepackage{array}
\usepackage{lipsum}
\usepackage{dsfont}
\usepackage{stmaryrd}

%\graphicspath{ {./diagrams/} }

\input{./latex-preamble/cm-preamble.tex}

\newcommand{\Tot}{\mb{Tot}}

\newtheorem*{lemma}{Lemma}

\begin{document}

{Weibel, \textit{An introduction to homological algebra}} \hfill {7-15-2026}

\bigskip\bigskip

\begin{center}

{\sc Chapter 1}

\end{center}

\begin{center}

  {Noah Parker} 
    
\end{center}

\bigskip

\begin{enumerate}
	\item[1.2.5.] {
      Let $C$ be a bounded complex with acylic rows or acyclic columns. 
      Show that $\Tot(C)$ is acyclic.
    }
\end{enumerate}

\begin{proof}
  Without loss of generality, we suppose that $C$ has acyclic rows, i.e. $\ker d^h=\im d^h$.
  It suffices to show that $\im(d^h+d^v)$ satisfies the universal property of $\ker(d^h+d^v)$.
  Now, let $f:E\to \Tot(C)_n$ be a morphism such that
  \[
    \begin{tikzcd}
      E\ar[r,"f"]\ar[rr,bend right=20,"0"']&\Tot (C)_n\ar[r,"d^h+d^v"]&\Tot(C)_{n-1}
    \end{tikzcd}
  \] 
  commutes.
  That is,
  \[
    d^hf+d^vf=0.
  \] 
  $C$ is bounded, so we can write
  \begin{align*}
    \Tot(C)_n =\bigoplus_{p+q=n}C_{pq}
              &=\prod_{p+q=n}C_{pq}
  \end{align*}
  as a product.
  In particular, we can write 
  \[
    f = \sum_{p+q=n}f_{pq}
  \] 
  for morphisms
  \[
    f_{pq}:E\to C_{pq}.
  \] 
  Then, for each $p+q=n$, the following diagram commutes and composes to 0.
  \[
    \begin{tikzcd}
      &C_{pq}\ar[dr,"d^h"]\\
      E\ar[ur,"f_{pq}"]\ar[dr,"f_{p+1,q-1}"']&&C_{p,q-1}\\%^{\oplus 2}\ar[r,"+"]&C_{p,q-1}\\
      &C_{p+1,q-1}\ar[ur,"d^v"']
    \end{tikzcd}
  \] 
  % Then, for each $p+q=n$, we have
  % \[
  %   d^hf_{pq}+d^vf_{p+1,q-1}=0.
  % \] 
  Suppose now that $f_{p+1,q-1}$ factors through $\im(d^h+d^v)$ as so:
  \[
    \begin{tikzcd}
      &C_{p+1,q}\ar[dr,"d^h"]\\
      E\ar[rr,"f_{p+1,q-1}"]\ar[ur,"g_{p+1,q}"]\ar[dr,"g_{p+2,q-1}"']&&\im(d^h+d^v)_{p+1,q-1}\ar[r,hook]&C_{p+1,q-1}\\
      &C_{p+2,q-1}\ar[ur,"d^v"']
    \end{tikzcd}
  \] 
  Then we have that
  \begin{align*}
    0 &= d^hf_{pq}+d^v f_{p+1,q-1} \\
      &= d^hf_{pq}+d^v(d^hg_{p+1,q} + d^vg_{p+2,q-1}) \\
      &= d^hf_{pq}+d^vd^hg_{p+1,q} \\
      &= d^h(f_{pq}-d^vg_{p+1,q}),
  \end{align*}
  whence $f_{pq}-d^vg_{p+1,q}\in \ker d^h = \im d^h$ gives us $g_{p,q+1}:E\to C_{p,q+1}$ such that $f_{pq}$ factors through $\im(d^h+d^v)_{pq}$ as follows:%, with $g_{p+1,q}:E\to C_{p+1,q}$ as before:
  \[
    \begin{tikzcd}
      &C_{p,q+1}\ar[dr,"d^h"]\\
      E\ar[rr,"f_{pq}"]\ar[ur,"g_{p,q+1}"]\ar[dr,"g_{p+1,q}"']&&\im(d^h+d^v)_{pq}\ar[r,hook]&C_{pq}\\
      &C_{p+1,q}\ar[ur,"d^v"']
    \end{tikzcd}
  \] 
  Note that this factorization holds trivially if $C_{pq}=0$.
  Then $C$ bounded supplies us with $q_0$ such that $C_{n-q,q}=0$ for all $q<q_0$, whence induction and the universal property of products give us a morphism
  \[
    g=\sum_{p+q=n+1}g_{pq}:E\to \Tot(C)_{n+1}
  \]
  such that
  \[
    \begin{tikzcd}
      E\ar[dr,"d^hg+d^vg"']\ar[r,"f"]&\Tot(C)_n\ar[r,"d^h+d^v"]&\Tot(C)_{n-1}\\
      &\im(d^h+d^v)\ar[u,hook]
    \end{tikzcd}
  \] 
  commutes.
  The map $d^hg+d^vg$ is unique since $\im(d^h+d^v)\inc \Tot(C)_n$ is a monomorphism, so $\im(d^h+d^v)\cong \ker(d^h+d^v)$ by uniqueness of the universal property of kernels.
  We conclude that $\Tot(C)_n$ is acyclic, as desired.
\end{proof}

\begin{enumerate}
	\item[1.2.7:] {
      Let $C$ be a complex. %, and define complexes $Z(C)$, $B(C)$, and $H(C)$ by $Z(C)_n:= \ker d_n$, $B(C)
      Show that there are exact sequences of complexes:
      \begin{equation}
        \begin{tikzcd}
          0\ar[r] & Z(C)\ar[r] & C\ar[r,"d"] & B(C)[-1]\ar[r] & 0
        \end{tikzcd}
      \end{equation}
      \begin{equation}
        \begin{tikzcd}
          0\ar[r] & H(C)\ar[r,"\phi"] & C/B(C)\ar[r,"d"] & Z(C)[-1]\ar[r,"\psi"] & H(C)[-1]\ar[r] & 0
        \end{tikzcd}
      \end{equation}
    }
\end{enumerate}

\begin{proof}
  Regard the differential as a map $d:C\to C[-1]$, so that $Z(C)=\ker d$ and $B(C)= \im d[1]$.
  Then $\im d=(\im d[1])[-1]=B(C)[-1]$ gives us well definition and exactness of (1).

  Consider now (2).
  We can regard $H(C)=Z(C)/B(C)=\cok (B(C)\inc Z(C))$.
  Now, we construct our map $\phi:H(C)\to C/B(C)$ by way of the following diagram:
  \[
    \begin{tikzcd}
      B(C)\ar[r,hook] & Z(C)\ar[r,hook]\ar[d,twoheadrightarrow] & C\ar[d,twoheadrightarrow] \\
                      & H(C) \ar[r, dotted, "\exists!\phi"] & C/B(C)
    \end{tikzcd}
  \] 
  Next, we will show that $d:C/B(C)\to Z(C)[-1]$ is well-defined.
  For this purpose, consider the following diagram:
  \[
    \begin{tikzcd}
      B(C)\ar[r,hook] & C\ar[r,"d"] & C[-1]\ar[r,"d{[-1]}"] & C[-2] \\
                      & C/B(C)\ar[u,twoheadrightarrow] \ar[ur,dotted]\ar[r,dotted,"d"'] & Z(C)[-1]\ar[u,hook]
    \end{tikzcd}
  \] 
  By abuse of notation, we will refer to each of these maps as $d$.
  The map $\psi:Z(C)[-1]\to H(C)[-1]$ is the cokernel of $B(C)[-1]\to Z(C)[-1]$, and thus epic.
  We may now proceed to showing exactness!

  We will show $\phi:H(C)\to C/B(C)$ is a monomorphism by showing it is the kernel of $d:C/B(C)\to Z(C)[-1]$.
  Before doing so, we first note that $\im d=\ker \psi$.
  Indeed, we have
  \[
    \begin{tikzcd}
      C\ar[r,twoheadrightarrow] & C/B(C)\ar[r,"d"] & Z(C)[-1]\ar[d,twoheadrightarrow] \\
                                && f
    \end{tikzcd}
  \] 
  \begin{align*}
    \ker \psi
      &= B(C)[-1] \\
      &= \im (d:C\to C[-1]) \\
      &= \im (d:)
  \end{align*}
\end{proof}

\end{document}
